% TOP_PROBLEM: {"id": 20001254, "title": "Prime-factor advice in power-residue quotients", "category_id": 15, "category": {"id": 15, "name": "computer_science", "display_name": "Computer Science", "description": "Computational complexity, algorithms, and theoretical CS.", "slug": "computer-science", "order_index": 15, "created_at": "2026-07-31T15:26:25.671Z"}, "status": "partially_solved"}
% FIRST_POSTED: null
% ENTRY_KIND: statement_only
% !TeX program = lualatex
\documentclass[11pt]{article}
\usepackage[margin=25mm]{geometry}
\usepackage{fontspec}
\setmainfont{Latin Modern Roman}
\usepackage{amsmath,amssymb,mathtools}
\usepackage{unicode-math}
\setmathfont{Latin Modern Math}
\usepackage{xurl,hyperref}
\hypersetup{colorlinks=true,urlcolor=blue}
\setcounter{secnumdepth}{0}
\setlength{\parindent}{0pt}
\setlength{\parskip}{5pt}
\setlength{\emergencystretch}{3em}
\newcommand{\R}{\mathbb R}
\newcommand{\N}{\mathbb N}
\newcommand{\Z}{\mathbb Z}
\newcommand{\Q}{\mathbb Q}
\newcommand{\C}{\mathbb C}
\newcommand{\D}{\mathbb D}
\newcommand{\F}{\mathbb F}
\newcommand{\sref}[2]{\hyperlink{source:#1:#2}{[S#2]}}
\newcommand{\cataloguescope}[1]{#1}
\begin{document}
% UnsolvedMath ID 20001254; source catalogue code AIM-COMPUTATION-0092
\setcounter{section}{20001253}
\section{Prime-factor advice in power-residue quotients}\label{20001254}
{\small UnsolvedMath ID 20001254 · Computer Science}
\subsection{Definitions and mathematical statement}

\paragraph{Shared notation.}$\mathbb N$, $\mathbb Z$, $\mathbb Q$, $\mathbb R$, and $\mathbb C$ denote the natural numbers, integers, rationals, reals, and complex numbers, respectively. All additional parameters and hypotheses are specified in the question.


\paragraph{Statement recorded in the supplied source.}
For a given prime $p$, how much auxiliary data suffices to reduce the solution of higher-degree equations modulo $p$ to an efficient computation, in analogy with a single quadratic nonresidue for quadratic equations? Does a deterministic subexponential-time algorithm find a quadratic nonresidue modulo $p$ without an unproved hypothesis?
\par\smallskip\textit{Formulation references:} \sref{20001254}{1}, \sref{20001254}{2}.
\subsection{Short English statement}
For a given prime $p$, how much auxiliary data suffices to reduce the solution of higher-degree equations modulo $p$ to an efficient computation, in analogy with a single quadratic nonresidue for quadratic equations? Does a deterministic subexponential-time algorithm find a quadratic nonresidue modulo $p$ without an unproved hypothesis?
\subsection{Sources}
\begin{itemize}
\item[\hypertarget{source:20001254:1}{S1}] ULAM.AI and the UnsolvedMath Contributors, UnsolvedMath: A Curated Collection of Open Mathematics Problems, record 20001254 (AIM-COMPUTATION-0092). Catalogue attribution under CC BY 4.0.
\url{https://huggingface.co/datasets/ulamai/UnsolvedMath}.
\item[\hypertarget{source:20001254:2}{S2}] Original problem formulation and mathematical context.
\url{https://aimath.org/WWN/primesinp/primesinp.pdf}.
\end{itemize}
\label{end:20001254}
\end{document}
